\section{Discussion and Threats to Validity}
\label{sec:discussion}

\paragraph{Why persona conditioning hurts.}
The base model, asked directly, holds a usable population-level prior on what gets clicked---
plausibly because clickability patterns are abundant in its pretraining data. Conditioning on a
specific persona (``you are a 68-year-old retiree\ldots'') reframes the task as first-person
role-play, which substitutes an idiosyncratic, stereotyped guess for that prior; averaging across
a ten-persona panel does not recover it, because each response is biased rather than merely noisy
---a caricature effect documented for LLM persona simulations~\citep{cheng2023compost}, where
persona conditioning can surface implicit bias and degrade performance on objective
tasks~\citep{gupta2024bias}. Seen another way, our no-persona baseline is itself a single
\emph{aggregate} persona (a ``typical reader''); the finding is then that \emph{disaggregating}
the audience into a demographic panel hurts an \emph{aggregate} prediction task---each
disaggregated draw adds role-play bias without adding signal the aggregate query lacks.
Empirically this is what we see: the panel's scores are no less spread within a package than the
baseline's (within-package SD \personaSpreadSD{} vs.\ \baseSpreadSD{}), so personas err by
\emph{systematic mis-ordering}, not indecision. In \nBaseRightPersonaWrong{} of \eaNPackages{}
reliable packages the baseline ranks the true winner first while the panel does not; these are
typically curiosity-gap or feel-good headlines (Appendix, Table~\ref{tab:error_examples}).
This is consistent with silicon-sampling working best for \emph{distributional} questions about
populations~\citep{argyleSiliconSampling} while underperforming for a \emph{point} prediction of
aggregate engagement, where the unconditioned prior is the better estimator. It also aligns with
independent reports that personas in prompts do not improve---and can
degrade---performance~\citep{zheng2024helpful}, and that the persona effect is small and scales
with the persona--outcome correlation~\citep{hu2024persona}. The latter predicts our
domain-dependence (Section~\ref{sec:results}): the baseline's advantage is large on controlled
copy A/B tests and shrinks where persona-relevant variance is lower.

\paragraph{Construct validity (RQ3).}
We elicit a \emph{click-intent} signal matched to the benchmark's engagement (click-through)
outcome, and claim relative ranking rather than absolute CTR~\citep{cronbach1955construct}. A
residual construct gap remains---headlines are editorial, and a \emph{stated} click-intent is not
identical to \emph{revealed} clicking---which bounds the achievable validity and plausibly
contributes to the small effect size. Tools that instead elicit a purchase- or conversion-intent
construct would face an even wider gap to a click outcome; matching the elicited construct to the
measured outcome is itself a design lesson.

\paragraph{Statistical-conclusion validity.}
An initial small-sample pilot overestimated the effect---roughly double the powered estimate---and
at the full sample it falls to $\tau=\kendallTau{}$, a cautionary instance of small-sample
optimism. We pre-registered
effect-size thresholds and report bootstrap CIs and seed-stability (SD $=\stabTauSD{}$ in $\tau$)
to make the weakness of the effect explicit. Top-1 accuracy is under-powered and reported as
secondary.

\paragraph{External validity and temporal mismatch.}
Evidence is from English-language editorial headlines circa 2013--2015, whereas the LLM was
trained on a later and broader slice of the web. This temporal gap cuts both ways: the model may
have absorbed post-2015 clickbait conventions that inflate the no-persona baseline, or it may
have ingested the Upworthy archive itself. We cannot fully rule out either, and flag it as a
threat; a contemporaneous benchmark would isolate it. Generalisation to marketing landing-page or
ad copy, other languages, and current audiences likewise requires further ground truth
(live deployment outcomes or a human panel), which we leave to future work.

\paragraph{Internal validity.}
Model, temperature, and all seeds are pinned and recorded per prediction; the persona panel is
derived from the benchmark audience~\citep{holcomb2013newsuse} rather than internal personas to
avoid authoring bias (Appendix~\ref{sec:appendix}). A model comparison found the result stable
across three model tiers (Table~\ref{tab:model_comparison}).

\paragraph{Design choices not exhausted.}
Our estimate fixes one elicitation prompt and one panel size. The aggregation rule, by contrast,
is no longer a free parameter: a sweep over six rules---including median, rank-fusion (Borda), and
pairwise majority (Condorcet)---leaves the conclusion intact (Table~\ref{tab:h5agg}, RQ3), and
locates the failure in genuine persona mis-ordering rather than the pooling step. Richer
psychographic personas, chain-of-thought elicitation, and larger panels could still shift the
effect; we report a single well-controlled configuration and leave a sensitivity sweep over those
remaining choices to future work. The failure is not specific to the Gemini family: re-running the
full protocol on a different family (OpenAI \xfModel{}) reproduces it on all \xfN{} reliable-winner
packages---the no-persona baseline again outranks the persona panel (Kendall $\tau$ \xfBaselineTau{}
vs.\ \xfPersonaTau{}; top-1 \xfBaselineTop{} vs.\ \xfPersonaTop{}), and the paired per-package
baseline$-$panel gap is significant in both metrics ($\tau$ gap \xfGapTau{}, 95\% CI
$[\xfGapTauCILow{}, \xfGapTauCIHigh{}]$; top-1 gap \xfGapTopPP{}pp, 95\% CI $[\xfGapTopPPCILow{},
\xfGapTopPPCIHigh{}]$; both exclude $0$). Effective persona adoption is therefore not simply an
emergent capability the tested tiers lack: a larger frontier model from an unrelated family fails the
same way. Whether \emph{any} model configuration makes persona conditioning beneficial for this
task remains open.
The companion baseline (Section~\ref{sec:results}) isolates the value of persona conditioning
against a no-persona zero-shot ranker.
